% Methodology (owner M). Every number below is read from a committed config,
% manifest or results file; the source is named in a comment beside it.

\section{Methodology}
\label{sec:method}

\subsection{Three systems}
\label{sec:systems}

We compare three detectors that share one architecture and differ only in what
they are trained on and which weights move.

\begin{itemize}
  \item \textbf{S1, English only.} Trained on ASVspoof 2019 LA~\cite{todisco2019asvspoof}
        and nothing else. It is the conventional detector whose transfer to
        code-mixed speech we measure.
  \item \textbf{S2, LoRA adaptation.} S1 with low-rank adapters~\cite{hu2022lora}
        trained on code-mixed Hindi--English data. The encoder weights of S1 stay
        frozen, so S2 can be read as S1 plus a small correction.
  \item \textbf{S3, native training.} The same architecture fine-tuned end to end on
        exactly the code-mixed manifests used by the matching S2 run, starting from
        the self-supervised checkpoint and never from S1. It shows what the
        code-mixed data alone supports.
\end{itemize}

S2 and S3 are each trained in four variants: on XTTS clones only or on XTTS clones
plus RVC conversions, and on clean audio or on channel-matched audio
(Sec.~\ref{sec:channel}). This gives eight code-mixed models and one English model.

\subsection{Data}
\label{sec:data}

\textbf{English.} ASVspoof 2019 LA~\cite{todisco2019asvspoof}: the official training
partition for S1, and the official evaluation partition of 71{,}237 clips over
attacks A07--A19 for English retention.

\textbf{Code-mixed real speech.} Hindi--English code-switched lecture speech from the
MUCS 2021 corpus~\cite{diwan2021mucs}. Speakers are divided into disjoint pools
before any synthesis: an adaptation pool of 10 speakers, used only by S2 and S3,
and an evaluation pool of 15 speakers, never seen in training.
% configs/data/splits.yaml; docs/datasheet.md

\textbf{Code-mixed synthetic speech.} Four attack sets are built from these pools
(Table~\ref{tab:attacks}). XTTS-v2~\cite{casanova2024xtts} zero-shot clones are made
for both pools, so at test time XTTS is a \emph{seen tool on unseen speakers}.
RVC~\cite{rvc2023} converts one MUCS speaker into another. Its conversions are split
by speaker so that neither the source nor the target speaker of any training
conversion appears in the RVC test half. Tortoise-TTS~\cite{betker2023tortoise}
clones of the evaluation pool form the held-out attack: by rule it never appears in
any training manifest, and an automated test fails the build if it does.
% tests/test_splits.py, tests/test_config_guard.py

\begin{table}[t]
\centering
\caption{Code-mixed attack sets. Usable clips passed the quality screen.}
\label{tab:attacks}
\small
\begin{tabular}{@{}llrl@{}}
\toprule
Set  & Tool      & Usable  & Role \\
\midrule
CM01 & XTTS-v2   & 3{,}998 & test, seen tool \\
CM02 & RVC       & 1{,}404 & split: train half, test half \\
CM03 & XTTS-v2   & 1{,}600 & adaptation only \\
CM04 & Tortoise  & 458     & test, unseen tool \\
\bottomrule
\end{tabular}
% docs/datasheet.md
\end{table}

\begin{table}[t]
\centering
\caption{Code-mixed training and test sets (clips).}
\label{tab:splits}
\small
\begin{tabular}{@{}lrrr@{}}
\toprule
Set & Real & Fake & Total \\
\midrule
Adaptation, XTTS only   & 874     & 1{,}280 & 2{,}154 \\
Adaptation, XTTS + RVC  & 1{,}217 & 1{,}637 & 2{,}854 \\
Adaptation dev          & 186     & 320     & 506 \\
\midrule
Eval pool (XTTS)        & 1{,}566 & 2{,}400 & 3{,}966 \\
CM04 (Tortoise)         & 1{,}566 & 458     & 2{,}024 \\
RVC test\textsuperscript{a} & 302     & 317     & 619 \\
\bottomrule
\multicolumn{4}{@{}l}{\footnotesize \textsuperscript{a}Real side: MUCS source clips of the held-out speakers.}
\end{tabular}
% data/manifests/codemix_adapt_*.csv, codemix_eval.csv; results/w10/*.json
\end{table}

\textbf{Preprocessing.} All audio is resampled to 16~kHz mono, trimmed by voice
activity detection and kept between 2 and 10~s. Every clip is level-normalised to
$-23$~dBFS RMS. This step matters: without it the XTTS clones arrive
peak-normalised while the lecture audio does not, and signal level alone separates
the classes (Sec.~\ref{sec:floor}).
% configs/data/channel_sim.yaml (preprocess block); docs/results/bundle_normalisation_v1.md

\subsection{Channel-matched condition}
\label{sec:channel}

The target application is a telephone call. We therefore render every clip through a
narrowband telephony chain: resample from 16 to 8~kHz, encode and decode with G.711
$\mu$-law, add white noise at 20~dB SNR and resample back to 16~kHz. Noise is seeded
per clip from a hash of its name. A single shared noise draw would add one identical
waveform to every clip, which a detector could learn in place of the speech. On the
evaluation render, energy above 4~kHz falls from 0.185\% to 0.0001\% and the measured
SNR is 18.8~dB. A model is always scored in the condition it was trained in.
% configs/data/channel_sim.yaml; docs/results/channel_matched_v1.md

\subsection{Model}
\label{sec:model}

The detector is a wav2vec~2.0 base encoder~\cite{baevski2020wav2vec} followed by
attentive statistics pooling~\cite{okabe2018attentive} and a two-layer classifier.
The encoder reads the raw 16~kHz waveform. Seven convolutional layers produce one
768-dimensional frame every 20~ms, each covering 25~ms of audio, and twelve
transformer layers contextualise the frames. Pooling computes an attention-weighted
mean and standard deviation over frames (attention width 128), giving a
1{,}536-dimensional vector. The classifier is Linear(1536, 256), ReLU, dropout 0.3 and
Linear(256, 2); its softmax output is read as $P(\text{genuine})$. The model has
95{,}454{,}083 parameters. In every system the convolutional feature extractor is
frozen, following common practice for self-supervised
anti-spoofing~\cite{tak2022wav2vec}.
% src/models/{encoder,pooling,classifier,detector}.py

\subsection{Training}
\label{sec:training}

All systems use AdamW with weight decay $10^{-5}$, 4~s training crops, mixed
precision and seed 1234. The remaining settings are in Table~\ref{tab:hparams}.
The checkpoint with the lowest EER on the development set is kept.

\textbf{S2.} Rank-8 adapters ($\alpha = 16$, dropout 0.05) wrap the query, key,
value and output projections of all twelve transformer layers. Only the adapters and
the classifier head are trained: 1{,}082{,}371 parameters, or 1.13\% of the model.
% src/models/lora.py; docs/results/gap_closure_v1.md

\textbf{S2 and S3} use class-weighted cross-entropy, since the adaptation sets are
not balanced (Table~\ref{tab:splits}).

\begin{table}[t]
\centering
\caption{Training settings.}
\label{tab:hparams}
\small
\begin{tabular}{@{}lccc@{}}
\toprule
                 & S1 & S2 (LoRA) & S3 (native) \\
\midrule
Initialised from & wav2vec 2.0 & S1 & wav2vec 2.0 \\
Trained weights  & transformer, head & adapters, head & transformer, head \\
Learning rate    & $10^{-5}$ & $10^{-4}$ & $10^{-5}$ \\
Batch size       & 4 & 8 & 8 \\
Epochs           & 3 & 5 & 10 \\
Training data    & ASVspoof LA & code-mixed & code-mixed \\
\bottomrule
\end{tabular}
% configs/train_stage1_run.yaml, train_lora_norm_*.yaml, train_s3_native_*.yaml
\end{table}

\subsection{Evaluation}
\label{sec:eval}

We report equal error rate (EER) and area under the ROC curve (AUC), scoring one
deterministic 4~s crop per clip, with 95\% bootstrap confidence intervals over clips.
Three code-mixed test sets probe three kinds of generalisation: the evaluation pool
(seen tool, unseen speakers), CM04 (unseen tool) and the RVC test half (speakers
unseen as either source or target). The ASVspoof 2019 LA evaluation set measures
what each code-mixed model retains of English.

\subsubsection{Shortcut floor}
\label{sec:floor}
A low EER on a newly built corpus is not evidence by itself. The real and fake halves
come from different pipelines, and pipeline provenance can separate them without any
knowledge of synthetic speech. Following the check proposed with
AffectDF~\cite{affectdf2026}, we fit a logistic regression on eight cheap signal
statistics (RMS mean and standard deviation, peak amplitude, clipping ratio, DC
offset, zero-crossing rate, spectral roll-off and high-frequency energy ratio) on the
adaptation pool, and score it on the speaker-disjoint test set. Its EER is that set's
\emph{shortcut floor}, measured separately for every test set and condition. We count
a detector's result as evidence only when the upper bound of its confidence interval
lies below the floor.
% src/data/lowlevel_cue.py; experiments/results/lowlevel_cue_check_*.json
